% Options for packages loaded elsewhere
\PassOptionsToPackage{unicode}{hyperref}
\PassOptionsToPackage{hyphens}{url}
\PassOptionsToPackage{dvipsnames,svgnames,x11names}{xcolor}
\documentclass[
  11pt,
]{article}
\usepackage{xcolor}
\usepackage[margin=2cm]{geometry}
\usepackage{amsmath,amssymb}
\setcounter{secnumdepth}{-\maxdimen} % remove section numbering
\usepackage{iftex}
\ifPDFTeX
  \usepackage[T1]{fontenc}
  \usepackage[utf8]{inputenc}
  \usepackage{textcomp} % provide euro and other symbols
\else % if luatex or xetex
  \usepackage{unicode-math} % this also loads fontspec
  \defaultfontfeatures{Scale=MatchLowercase}
  \defaultfontfeatures[\rmfamily]{Ligatures=TeX,Scale=1}
\fi
\usepackage{lmodern}
\ifPDFTeX\else
  % xetex/luatex font selection
\fi
% Use upquote if available, for straight quotes in verbatim environments
\IfFileExists{upquote.sty}{\usepackage{upquote}}{}
\IfFileExists{microtype.sty}{% use microtype if available
  \usepackage[]{microtype}
  \UseMicrotypeSet[protrusion]{basicmath} % disable protrusion for tt fonts
}{}
\makeatletter
\@ifundefined{KOMAClassName}{% if non-KOMA class
  \IfFileExists{parskip.sty}{%
    \usepackage{parskip}
  }{% else
    \setlength{\parindent}{0pt}
    \setlength{\parskip}{6pt plus 2pt minus 1pt}}
}{% if KOMA class
  \KOMAoptions{parskip=half}}
\makeatother
\usepackage{listings}
\newcommand{\passthrough}[1]{#1}
\lstset{defaultdialect=[5.3]Lua}
\lstset{defaultdialect=[x86masm]Assembler}
\lstset{
  breaklines=true,
  breakatwhitespace=false,
  columns=fullflexible,
  basicstyle=\ttfamily\small,
  keepspaces=true,
  showstringspaces=false,
  frame=single,
  postbreak=\mbox{\textcolor{gray}{$\hookrightarrow$}\space}
}
\lstset{
  breaklines=true,
  breakatwhitespace=false,
  columns=fullflexible,
  basicstyle=\ttfamily\small,
  keepspaces=true,
  showstringspaces=false,
  frame=single,
  postbreak=\mbox{\textcolor{gray}{$\hookrightarrow$}\space}
}
\usepackage{longtable,booktabs,array}
\usepackage{caption}
\captionsetup[table]{skip=6pt}
\newcounter{none} % for unnumbered tables
\usepackage{calc} % for calculating minipage widths
% Correct order of tables after \paragraph or \subparagraph
\usepackage{etoolbox}
\makeatletter
\patchcmd\longtable{\par}{\if@noskipsec\mbox{}\fi\par}{}{}
\makeatother
\setlength{\emergencystretch}{8em} % prevent overfull lines
\providecommand{\tightlist}{%
  \setlength{\itemsep}{0pt}\setlength{\parskip}{0pt}}
\usepackage{bookmark}
\IfFileExists{xurl.sty}{\usepackage{xurl}}{} % add URL line breaks if available
\urlstyle{same}
% fallback for those not using the hyperref driver hyperxmp:
\makeatletter
\@ifundefined{xmpquote}{\newcommand{\xmpquote}[1]{#1}}{}
\makeatother
\hypersetup{
  pdftitle={Destructuring Field Guide},
  pdfauthor={Abdullah Olcan Askar},
  colorlinks=true,
  linkcolor={blue},
  filecolor={Maroon},
  citecolor={Blue},
  urlcolor={blue},
  pdfcreator={LaTeX via pandoc}}

\title{Destructuring Field Guide}
\author{Abdullah Olcan Askar\\Student number: 2024502180\\MIS 233, Fall 2026}
\date{November 2026}

\begin{document}
\raggedright
\raggedright
\maketitle
\thispagestyle{empty}
\newpage

{
\setcounter{tocdepth}{3}
\tableofcontents
}
\newpage

\section{Repository}\label{repository}

https://github.com/olcanaskar/mis233-assignment4

The repository is public.

\section{Run}\label{run}

\begin{lstlisting}
npm install
npm run all
npm run check
\end{lstlisting}

\texttt{output.txt} is the real output of
\texttt{npm run all}.

\section{Data}\label{data}

My last ten films on Letterboxd
(\texttt{letterboxd.com/olcan}), newest first:

Tokyo Godfathers (2003, 4), Primetime (2026, 3.5), Memories of Murder
(2003, 4.5), The Worst Person in the World (2021, 4), The Lobster (2015,
3.5), Coyote vs.~Acme (2026, 3.5), Happy Together (1997, 3), The Working
Class Goes to Heaven (1971, 3), Punch-Drunk Love (2002, 3.5), Waves
(2019, 4).

The list is in \texttt{js/data.mjs} and
\texttt{ts/data.ts}. The TypeScript data only has
\texttt{id}, \texttt{title},
\texttt{year}, \texttt{director.name},
\texttt{genres} and \texttt{rating}.

\section{Field guide}\label{field-guide}

All outputs are from \texttt{output.txt}.

\subsection{\texorpdfstring{Form 1: values by position,
\texttt{const\ {[}x,\ y{]}\ =\ f()}}{Form 1: values by position, const {[}x, y{]} = f()}}\label{form-1-values-by-position-const-x-y-f}

\textbf{Syntax:}
\texttt{const [oldest, newest] = yearRange(films)}

\textbf{Example} (\texttt{js/arrays.mjs}):
\texttt{yearRange} returns
\texttt{[Math.min(...years), Math.max(...years)]}, then
I destructure it.

\textbf{Output:} \texttt{form 1: 1971 2026}

\textbf{Use it when:} a function returns several values in a fixed
order.

\textbf{Watch out:} order is the only link. A missing item is
\texttt{undefined}. In TS I needed
\texttt{: [number, number]}; otherwise the positions
are \texttt{number \textbar{} undefined}.

\subsection{Form 2: skip an item and collect the
rest}\label{form-2-skip-an-item-and-collect-the-rest}

\textbf{Syntax:}
\texttt{const [, second, ...others] = ranked}

\textbf{Example} (\texttt{js/arrays.mjs}): I sort by
rating descending, skip the first film, and take the second plus the
rest.

\textbf{Output:}
\texttt{form 2: Tokyo Godfathers + 8 more}

\textbf{Use it when:} the first item is not important but the rest are.

\textbf{Watch out:} the comma matters;
\texttt{[second, ...others]} would take the first film.
In TS \texttt{second} is
\texttt{Film \textbar{} undefined}, so I used
\texttt{second?.title}.

\subsection{Form 3: swap two variables}\label{form-3-swap-two-variables}

\textbf{Syntax:} \texttt{[a, b] = [b, a]}

\textbf{Example} (\texttt{js/arrays.mjs}):
\texttt{let firstPick = ...; let secondPick = ...; [firstPick, secondPick] = [secondPick, firstPick]}

\textbf{Output:}
\texttt{form 3: Primetime \textless{}-\textgreater{} Tokyo Godfathers}

\textbf{Use it when:} two values need to change places.

\textbf{Watch out:} use \texttt{let}, not
\texttt{const}. The swap line starts with
\texttt{;} because the previous line ends with
\texttt{]}. In TS I checked that both films exist
first.

\subsection{Form 4: default values in an
array}\label{form-4-default-values-in-an-array}

\textbf{Syntax:}
\texttt{const [firstYear = 2003, secondYear = 2026] = emptyYears}

\textbf{Example} (\texttt{js/arrays.mjs}):
\texttt{const emptyYears = []} then the line above.

\textbf{Output:} \texttt{form 4: 2003 2026}

\textbf{Use it when:} an item may be missing and there is a fallback.

\textbf{Watch out:} the default only runs for
\texttt{undefined}, not \texttt{null}.

\subsection{Form 5: properties by name}\label{form-5-properties-by-name}

\textbf{Syntax:}
\texttt{const \{ title, year \} = filmAt(0)}

\textbf{Example} (\texttt{js/objects.mjs}):
\texttt{const \{ title, year \} = filmAt(0)}

\textbf{Output:} \texttt{form 5: Tokyo Godfathers 2003}

\textbf{Use it when:} the property names are the names you want.

\textbf{Watch out:} a typo gives \texttt{undefined} in
JS. If the value is \texttt{undefined} or
\texttt{null}, the line throws.

\subsection{Form 6: rename while you take
it}\label{form-6-rename-while-you-take-it}

\textbf{Syntax:}
\texttt{const \{ title: filmTitle, year: releaseYear \} = filmAt(1)}

\textbf{Example} (\texttt{js/objects.mjs}):
\texttt{const \{ title: filmTitle, year: releaseYear \} = filmAt(1)}

\textbf{Output:} \texttt{form 6: Primetime 2026}

\textbf{Use it when:} the property name clashes or is not clear.

\textbf{Watch out:} left of the colon is the property; right is the new
variable. After this line there is no \texttt{title} or
\texttt{year} variable. In TS,
\texttt{\{ title: string \}} inside a pattern is also a
rename, not a type.

\subsection{Form 7: default values in an
object}\label{form-7-default-values-in-an-object}

\textbf{Syntax:}
\texttt{const \{ rating = 0 \} = unrated}

\textbf{Example} (\texttt{js/objects.mjs}):
\texttt{const unrated = \{ ...filmAt(0), rating: undefined \}}

\textbf{Output:} \texttt{form 7: 0}

\textbf{Use it when:} a field is optional.

\textbf{Watch out:} \texttt{null} does not trigger the
default.

\subsection{Form 8: rest in an object}\label{form-8-rest-in-an-object}

\textbf{Syntax:}
\texttt{const \{ id, ...filmWithoutId \} = filmAt(3)}

\textbf{Example} (\texttt{js/objects.mjs}):
\texttt{const \{ id, ...filmWithoutId \} = filmAt(3)}

\textbf{Output:}
\texttt{form 8: the-worst-person-in-the-world title,year,director,genres,rating}

\textbf{Use it when:} one property must be removed.

\textbf{Watch out:} rest must be last, and the copy is shallow. Nested
objects are still shared.

\subsection{Form 9: nested patterns}\label{form-9-nested-patterns}

\textbf{Syntax:}
\texttt{const \{ director: \{ name: directorName \}, genres: [mainGenre] \} = filmAt(2)}

\textbf{Example} (\texttt{js/objects.mjs}): the line
above.

\textbf{Output:} \texttt{form 9: Bong Joon-ho crime}

\textbf{Use it when:} you only need a deep piece of an object.

\textbf{Watch out:} every level must exist.
\texttt{mainGenre} can be
\texttt{undefined} if the array is empty; TS types it
as \texttt{string \textbar{} undefined}.

\subsection{Form 10: a computed key}\label{form-10-a-computed-key}

\textbf{Syntax:}
\texttt{const key = 'year'; const \{ [key]: value \} = filmAt(9)}

\textbf{Example} (\texttt{js/objects.mjs}): the line
above.

\textbf{Output:} \texttt{form 10: year 2019}

\textbf{Use it when:} the property name is chosen at runtime.

\textbf{Watch out:} the brackets are required. In TS a plain
\texttt{string} key gives
\texttt{TS2537}; use a literal or
\texttt{keyof Film}.

\subsection{Form 11: patterns in
parameters}\label{form-11-patterns-in-parameters}

\textbf{Syntax:}
\texttt{function createFilm(\{ title = 'Untitled', year = 0, rating = 0 \} = \{\})}

\textbf{Example} (\texttt{js/params.mjs}):
\texttt{createFilm(\{ title: 'Waves', year: 2019 \})}
and \texttt{createFilm()}.

\textbf{Output:}
\texttt{form 11: \{ title: 'Waves', year: 2019, rating: 0 \}}
and
\texttt{form 11 empty: \{ title: 'Untitled', year: 0, rating: 0 \}}

\textbf{Use it when:} a function takes an options object with defaults.

\textbf{Watch out:} \texttt{= \{\}} is a separate
default. In TS the type goes after the pattern; if the type has required
fields, \texttt{= \{\}} does not compile.
\texttt{null} throws.

\textbf{Form 11b:}
\texttt{describePair([first, second])} with two films
gives \texttt{form 11b: Tokyo Godfathers / Primetime}.

\subsection{Form 12: patterns in loops}\label{form-12-patterns-in-loops}

\textbf{Syntax:}
\texttt{for (const [id, film] of Object.entries(byId))}

\textbf{Example} (\texttt{js/params.mjs}): I build
\texttt{byId} with
\texttt{Object.fromEntries}, then loop. I also loop
over a \texttt{Map} with
\texttt{[title, year]}.

\textbf{Output:}
\texttt{form 12: tokyo-godfathers =\textgreater{} Tokyo Godfathers}
and \texttt{form 12 map: Tokyo Godfathers 2003} (full
list in \texttt{output.txt}).

\textbf{Use it when:} you need both key and value in a loop.

\textbf{Watch out:} \texttt{Object.entries} gives a
string key. In TS a \texttt{Map} needs the tuple type
\texttt{[string, number]}.
\texttt{for...in} gives keys only.

\section{Pitfalls}\label{pitfalls}

\subsection{\texorpdfstring{1. Destructuring \texttt{undefined} or
\texttt{null}}{1. Destructuring undefined or null}}\label{destructuring-undefined-or-null}

\texttt{const \{ title \} = undefined} gives
\texttt{TypeError: Cannot destructure property 'title' of 'undefined' as it is undefined.}
\texttt{null} gives the same error. \textbf{Why:} they
have no properties.

\subsection{\texorpdfstring{2. Default is only for
\texttt{undefined}}{2. Default is only for undefined}}\label{default-is-only-for-undefined}

\texttt{const \{ rating = 5 \} = \{ rating: null \}}
prints \texttt{null}. \textbf{Why:}
\texttt{null} is a real value, not a missing one.

\subsection{\texorpdfstring{3. Assignment without
\texttt{const}/\texttt{let} needs
parentheses}{3. Assignment without const/let needs parentheses}}\label{assignment-without-constlet-needs-parentheses}

\texttt{\{ title \} = film} gives
\texttt{SyntaxError: Unexpected token '='};
\texttt{(\{ title \} = film)} works. \textbf{Why:} a
line starting with \texttt{\{} is read as a block.

\subsection{4. Destructuring is a shallow
copy}\label{destructuring-is-a-shallow-copy}

\texttt{const \{ director \} = film; director.name = 'X'}
changes the original film too. \textbf{Why:} the nested object is shared
by reference.

\subsection{5. Array order matters, object order does
not}\label{array-order-matters-object-order-does-not}

\texttt{[a, b] = ['A', 'B']} changes if you swap the
array. \texttt{\{ first, second \}} from
\texttt{\{ second: 'B', first: 'A' \}} still gives the
same values. \textbf{Why:} arrays use position, objects use names.

\section{TypeScript topics}\label{typescript-topics}

\subsection{\texorpdfstring{T1: \texttt{\{\ a:\ string\ \}} is a rename,
not a
type}{T1: \{ a: string \} is a rename, not a type}}\label{t1-a-string-is-a-rename-not-a-type}

\texttt{const \{ title: string \} = filmAt(0)} prints
\texttt{T1 rename: Tokyo Godfathers}. It creates a
variable named \texttt{string}. The type goes after the
pattern:
\texttt{const \{ title: filmTitle \}: Film = filmAt(0)}.

\subsection{\texorpdfstring{T2: tuple and
\texttt{as\ const}}{T2: tuple and as const}}\label{t2-tuple-and-as-const}

\texttt{const pair: [number, string] = [2003, 'Tokyo Godfathers']}
prints \texttt{T2 tuple: 2003 Tokyo Godfathers}.
\texttt{const years = [2003, 2026, 1997] as const}
prints \texttt{T2 as const: 2003 2026}. Without the
tuple type the positions are
\texttt{number \textbar{} undefined}.

\subsection{\texorpdfstring{T3:
\texttt{noUncheckedIndexedAccess}}{T3: noUncheckedIndexedAccess}}\label{t3-nouncheckedindexedaccess}

\texttt{const [firstFilm] = films} makes
\texttt{firstFilm} a
\texttt{Film \textbar{} undefined}.
\texttt{firstFilm.title} alone gives
\texttt{error TS18048: 'firstFilm' is possibly 'undefined'.}
I used \texttt{?.}, an \texttt{if}
check, or a default.

\subsection{T4: optional fields and
defaults}\label{t4-optional-fields-and-defaults}

\texttt{rating?: number} is
\texttt{number \textbar{} undefined}.
\texttt{const \{ rating = 0 \}} makes it
\texttt{number}.
\texttt{rating?: number \textbar{} null} with a default stays
\texttt{number \textbar{} null}. Outputs:
\texttt{T4 optional: 4.5},
\texttt{T4 empty: undefined 0},
\texttt{T4 null: null}.

\subsection{T5: a compile error caused by
destructuring}\label{t5-a-compile-error-caused-by-destructuring}

\texttt{const \{ rating \}: \{ rating: number \} = filmAt(0)}
fails because \texttt{rating} is optional:

\begin{lstlisting}
error TS2322: Type 'Film' is not assignable to type '{ rating: number; }'.
  Types of property 'rating' are incompatible.
    Type 'number | undefined' is not assignable to type 'number'.
      Type 'undefined' is not assignable to type 'number'.
\end{lstlisting}

The fixed line is
\texttt{const \{ rating: safeRating = 0 \}: \{ rating?: number \} = filmAt(0)}.
Output: \texttt{T5 suppressed: 4} and
\texttt{T5 fixed: 4}. I used
\texttt{@ts-expect-error} only on the failing line.

\section{Props and children}\label{props-and-children}

A component gets one object, \texttt{props}. Content
between the tags is the \texttt{children} field.

\subsection{P1: props is one object}\label{p1-props-is-one-object}

\texttt{function FilmCardA(props)} and
\texttt{function FilmCardB(\{ film: \{ title \}, children \})}
render the same thing:

\begin{lstlisting}
P1 A: <article><h2>Tokyo Godfathers</h2>same text</article>
P1 B: <article><h2>Tokyo Godfathers</h2>same text</article>
\end{lstlisting}

The second one uses form 11, form 9 and form 5.

\subsection{P2: children is only a
prop}\label{p2-children-is-only-a-prop}

\texttt{\textless{}FilmCardB film=\{filmAt(1)\}\textgreater{}Primetime note\textless{}/FilmCardB\textgreater{}}
and
\texttt{\textless{}FilmCardB film=\{filmAt(1)\} children="Primetime note" /\textgreater{}}
both give:

\begin{lstlisting}
P2 tag: <article><h2>Primetime</h2>Primetime note</article>
P2 prop: <article><h2>Primetime</h2>Primetime note</article>
\end{lstlisting}

In JS, giving both makes the tag content win. TS rejects it with
\texttt{TS2710: 'children' are specified twice}.

\subsection{P3: what children can be}\label{p3-what-children-can-be}

{\def\LTcaptype{none} % do not increment counter
\begin{longtable}[]{@{}>{\raggedright\arraybackslash}p{0.28\linewidth}>{\raggedright\arraybackslash}p{0.66\linewidth}@{}}
\toprule\noalign{}
children & result \\
\midrule\noalign{}
\endhead
\bottomrule\noalign{}
\endlastfoot
string, number, element, array & rendered \\
\texttt{undefined}, \texttt{null},
\texttt{false} & nothing \\
function & warning, nothing \\
\end{longtable}
}

For example:
\texttt{P3 string: "\textless{}article\textgreater{}\textless{}h2\textgreater{}Tokyo Godfathers\textless{}/h2\textgreater{}a review\textless{}/article\textgreater{}"}
and
\texttt{P3 undefined: "\textless{}article\textgreater{}\textless{}h2\textgreater{}Tokyo Godfathers\textless{}/h2\textgreater{}\textless{}/article\textgreater{}"}.
A function produces
\texttt{Functions are not valid as a React child.}
\texttt{ReactNode} does not include functions.

\subsection{P4: children vs other
props}\label{p4-children-vs-other-props}

\texttt{(\{ film, children, ...rest \})} removes
\texttt{film} and \texttt{children};
\texttt{rest} has \texttt{className}
and \texttt{id}.
\texttt{\textless{}section \{...rest\}\textgreater{}} gives:

\begin{lstlisting}
P4 rest: <section class="card" id="tokyo"><h2>Tokyo Godfathers</h2>text</section>
\end{lstlisting}

Spread order decides the winner:
\texttt{\textless{}section \{...rest\} className="explicit"\textgreater{}}
keeps \texttt{explicit};
\texttt{\textless{}section className="explicit" \{...rest\}\textgreater{}}
keeps the rest value.

\subsection{P5: props or children?}\label{p5-props-or-children}

\texttt{SlotCard} takes
\texttt{header} and \texttt{footer};
\texttt{ChildrenCard} takes children only. Outputs:

\begin{lstlisting}
P5 slots: <article><header><span>Header</span></header><h2>The Worst Person in the World</h2><footer><span>Footer</span></footer></article>
P5 children: <article><h2>The Worst Person in the World</h2><span>Header</span><span>Footer</span></article>
\end{lstlisting}

Use slots when the layout is fixed and each part has a job. Use children
when the caller should decide the content and order.

\subsection{\texorpdfstring{P6: defaults, \texttt{undefined} and
\texttt{null}}{P6: defaults, undefined and null}}\label{p6-defaults-undefined-and-null}

\begin{lstlisting}
P6 default: <span>0</span>
P6 undefined: <span>0</span>
P6 null: <span>no rating</span>
\end{lstlisting}

\texttt{undefined} triggers the default;
\texttt{null} does not. \texttt{??}
shows null as \texttt{no rating}.

\subsection{P7: types for props and
children}\label{p7-types-for-props-and-children}

\texttt{children?: ReactNode} accepts strings, numbers,
elements, arrays, \texttt{null},
\texttt{undefined} and \texttt{false};
it rejects functions.
\texttt{PropsWithChildren\textless{}\{ film: Film \}\textgreater{}} adds
optional children. \texttt{ComponentProps\textless{}'button'\textgreater{}}
gives all button props (\texttt{type},
\texttt{disabled}, \texttt{onClick},
\texttt{children}).

Outputs:

\begin{lstlisting}
P7 PropsWithChildren: <article><h2>Tokyo Godfathers</h2>child</article>
P7 ComponentProps button: <button type="button" disabled="">Save</button>
P7 key probe: <span>undefined</span>
\end{lstlisting}

The compiler rejects:

\begin{lstlisting}
<FilmCard film="Tokyo Godfathers" />  -> Type 'string' is not assignable to type 'Film'.
<FilmCard film={filmAt(0)} title="x" /> -> Property 'title' does not exist.
<FilmCard film={filmAt(0)}>{() => 'x'}</FilmCard> -> Type '() => string' is not assignable to type 'ReactNode'.
<FancyButton variant="ghost" /> -> Property 'variant' does not exist.
\end{lstlisting}

\texttt{key} is not a prop:
\texttt{props.key} is
\texttt{undefined} and React warns about it.

\section{JavaScript vs TypeScript}\label{javascript-vs-typescript}

{\def\LTcaptype{none} % do not increment counter
\begin{longtable}[]{@{}
  >{\raggedright\arraybackslash}p{(\linewidth - 6\tabcolsep) * \real{0.2500}}
  >{\raggedright\arraybackslash}p{(\linewidth - 6\tabcolsep) * \real{0.2500}}
  >{\raggedright\arraybackslash}p{(\linewidth - 6\tabcolsep) * \real{0.2500}}
  >{\raggedright\arraybackslash}p{(\linewidth - 6\tabcolsep) * \real{0.2500}}@{}}
\toprule\noalign{}
\begin{minipage}[b]{\linewidth}\raggedright
Form
\end{minipage} & \begin{minipage}[b]{\linewidth}\raggedright
Same in both
\end{minipage} & \begin{minipage}[b]{\linewidth}\raggedright
TypeScript adds
\end{minipage} & \begin{minipage}[b]{\linewidth}\raggedright
Error if wrong
\end{minipage} \\
\midrule\noalign{}
\endhead
\bottomrule\noalign{}
\endlastfoot
1. array by position & \texttt{const [a, b] = f()} &
\texttt{: [number, number]} &
\texttt{TS18048} if a position is used as a number. \\
2. skip and rest &
\texttt{const [, second, ...others] = ranked} & nothing
& \texttt{second.title} gives
\texttt{TS18048}. \\
3. swap & \texttt{[a, b] = [b, a]} &
\texttt{let} plus a check &
\texttt{films[0]} gives
\texttt{TS18048} if used directly. \\
6. rename & \texttt{const \{ a: alpha \} = t} & nothing
& \texttt{\{ a: string \}} is a rename, not a type. \\
7. object default & \texttt{const \{ a = 1 \} = t} &
\texttt{a?: number} & Required
\texttt{\{ a: number \}} gives
\texttt{TS2322}. \\
9. nested pattern & \texttt{const \{ p: \{ x \} \} = t}
& nothing & Array element is
\texttt{T \textbar{} undefined}. \\
10. computed key &
\texttt{const \{ [key]: value \} = t} & literal or
\texttt{keyof Film} & Plain
\texttt{string} gives
\texttt{TS2537}. \\
11. parameter pattern &
\texttt{function f(\{ a = 1 \} = \{\})} & type after
the pattern & Required fields make \texttt{= \{\}}
invalid. \\
12. loop pattern &
\texttt{for (const [k, v] of Object.entries(o))} &
tuple type for \texttt{Map} &
\texttt{Map} infers a union without it. \\
props and children & \texttt{\textless{}Card\textgreater{}text\textless{}/Card\textgreater{}} &
\texttt{children?: ReactNode} & Function child gives
\texttt{TS2322}; duplicate children give
\texttt{TS2710}. \\
\end{longtable}
}

\section{Bonus 2: Depth}\label{bonus-2-depth}

The handout asks what an array pattern really calls on the value. The
answer is the iterator protocol, not array indexes. I put the
experiments in \texttt{js/depth.mjs} and added that
file to \texttt{npm run js}.

\textbf{Output from \texttt{node js/depth.mjs}:}

\begin{lstlisting}
depth string: F i
depth set: drama comedy
depth map: Tokyo Godfathers 2003 | Primetime 2026
depth generator: Tokyo Godfathers 2003
depth custom: Memories of Murder 2003
depth rest: 1 2,3
depth object: TypeError - plainObject is not iterable
\end{lstlisting}

What each line shows:

\begin{itemize}
\tightlist
\item
  A string is iterable character by character, so
  \texttt{'Film'} gives \texttt{F} and
  \texttt{i}.
\item
  A \texttt{Set} is iterable; duplicates are removed
  before iteration, so the first two values are
  \texttt{drama} and \texttt{comedy}.
\item
  A \texttt{Map} yields
  \texttt{[key, value]} pairs, so the nested pattern
  \texttt{[[firstTitle, firstYear], [secondTitle, secondYear]]}
  works.
\item
  A generator is iterable without being an array; the pattern calls
  \texttt{next()} and takes the yielded values.
\item
  The custom object in \texttt{js/depth.mjs} has only
  \texttt{[Symbol.iterator]} and
  \texttt{next()}: no \texttt{length},
  no \texttt{[0]}. The array pattern still works.
\item
  \texttt{...rest} keeps pulling values from the
  iterator until it is done.
\item
  A plain object has no \texttt{Symbol.iterator}, so
  the pattern throws \texttt{TypeError}.
\end{itemize}

What I learned: \texttt{const [a, b] = value} does not
mean ``read \texttt{value[0]} and
\texttt{value[1]}''. It means ``get an iterator from
\texttt{value}, call \texttt{next()}
twice, and bind the results''. That is why strings,
\texttt{Set}, \texttt{Map}, generators
and custom iterables work, and plain objects do not. The same protocol
powers \texttt{for...of}, spread and
\texttt{Array.from}.

\section{Bonus 3: Real code}\label{bonus-3-real-code}

I used the TypeScript compiler installed in
\texttt{node\textbackslash{}\_modules} (version 5.8.3) and MUI's
open-source Button component. The TypeScript line numbers are from the
installed file; the MUI link is pinned to commit
\texttt{9f8fe58}.

\subsection{\texorpdfstring{1.
\texttt{node\_modules/typescript/lib/typescript.js:4872} --- skip and
defaults (forms 2 and
4)}{1. node\_modules/typescript/lib/typescript.js:4872 --- skip and defaults (forms 2 and 4)}}\label{node_modulestypescriptlibtypescript.js4872-skip-and-defaults-forms-2-and-4}

\begin{lstlisting}
const [, major, minor = "0", patch = "0", prerelease = "", build2 = ""] = match;
\end{lstlisting}

\texttt{match} comes from
\texttt{versionRegExp.exec(text)}. Index 0 is the whole
match, so the leading comma skips it. Without destructuring:

\begin{lstlisting}
const major = match[1];
const minor = match[2] === undefined ? "0" : match[2];
const patch = match[3] === undefined ? "0" : match[3];
const prerelease = match[4] === undefined ? "" : match[4];
const build2 = match[5] === undefined ? "" : match[5];
\end{lstlisting}

The destructured line is shorter and keeps each default next to its
name.

\subsection{\texorpdfstring{2.
\texttt{node\_modules/typescript/lib/typescript.js:4830} --- rename
(forms 5 and
6)}{2. node\_modules/typescript/lib/typescript.js:4830 --- rename (forms 5 and 6)}}\label{node_modulestypescriptlibtypescript.js4830-rename-forms-5-and-6}

\begin{lstlisting}
const { major, minor, patch, prerelease, build: build2 } = result;
\end{lstlisting}

Without destructuring:

\begin{lstlisting}
const major = result.major;
const minor = result.minor;
const patch = result.patch;
const prerelease = result.prerelease;
const build2 = result.build;
\end{lstlisting}

\texttt{build} is renamed to
\texttt{build2}, probably to avoid a name clash. This
is form 6 in real code.

\subsection{\texorpdfstring{3.
\texttt{node\_modules/typescript/lib/typescript.js:188380} --- rest
(form
8)}{3. node\_modules/typescript/lib/typescript.js:188380 --- rest (form 8)}}\label{node_modulestypescriptlibtypescript.js188380-rest-form-8}

\begin{lstlisting}
const { lazyConfiguredProjectsFromExternalProject: _, ...userPreferences } = preferences;
\end{lstlisting}

Without destructuring:

\begin{lstlisting}
const userPreferences = { ...preferences };
delete userPreferences.lazyConfiguredProjectsFromExternalProject;
\end{lstlisting}

The rest pattern removes one internal field and keeps the rest in a new
object. The \texttt{\textbackslash{}\_} name says ``I do not use this
value''.

\subsection{\texorpdfstring{4.
\texttt{node\_modules/typescript/lib/typescript.js:136407} --- loop
(form
12)}{4. node\_modules/typescript/lib/typescript.js:136407 --- loop (form 12)}}\label{node_modulestypescriptlibtypescript.js136407-loop-form-12}

\begin{lstlisting}
for (const [key, value] of lineCounts.entries()) {
  reportCountStatistic("Lines of " + key, value);
}
\end{lstlisting}

Without destructuring:

\begin{lstlisting}
for (const entry of lineCounts.entries()) {
  const key = entry[0];
  const value = entry[1];
  reportCountStatistic("Lines of " + key, value);
}
\end{lstlisting}

The destructured version names the two parts immediately.

\subsection{\texorpdfstring{5. MUI Button, \texttt{Button.js:507-527}
--- React props and children (forms 5, 6, 7, 8 and
11)}{5. MUI Button, Button.js:507-527 --- React props and children (forms 5, 6, 7, 8 and 11)}}\label{mui-button-button.js507-527-react-props-and-children-forms-5-6-7-8-and-11}

Source:
\href{https://github.com/mui/material-ui/blob/9f8fe586afb3d81f4108a96f3b943a0990ee3f08/packages/mui-material/src/Button/Button.js\#L507-L527}{\texttt{packages/mui-material/src/Button/Button.js}}

\begin{lstlisting}
const {
  children,
  color = 'primary',
  component = 'button',
  // ... about 16 more props ...
  variant = 'text',
  ...other
} = props;
\end{lstlisting}

This is a real React component. \texttt{props} is the
one object the component receives. \texttt{children} is
one field of that object, not a separate argument.
\texttt{color}, \texttt{component},
\texttt{variant} and the rest are props.
\texttt{...other} keeps the props that were not named;
later the component spreads them onto
\texttt{\textless{}ButtonRoot \{...other\}\textgreater{}}, while
\texttt{children} is rendered explicitly at line 608.

Without destructuring it would start like this:

\begin{lstlisting}
const children = props.children;
const color = props.color === undefined ? 'primary' : props.color;
const component = props.component === undefined ? 'button' : props.component;
// ... about 16 more lines ...
const other = { ...props };
delete other.children;
delete other.color;
// ...
\end{lstlisting}

\textbf{Where it is worse:} this is my example of destructuring making
code harder to read. The block has about 20 names, several defaults, a
rest element, and a \texttt{children} field mixed in. A
reader has to scan the whole block to know what is forwarded. MUI
probably does it to apply defaults and filter props, but for a component
this large I would prefer a few targeted reads and an explicit
\texttt{other} object.

\section{My own log}\label{my-own-log}

\begin{enumerate}
\def\labelenumi{\arabic{enumi}.}
\item
  \textbf{TS2710.} I put both
  \texttt{children="from prop"} and tag content in
  \texttt{ts/props.tsx}.
  \texttt{npm run check} said:
  \texttt{error TS2710: 'children' are specified twice.}
  I removed it from TS; JS still runs it and the tag content wins.
\item
  \textbf{Function child.} I passed
  \texttt{() =\textgreater{} 'ignored'} as children. React printed
  \texttt{Functions are not valid as a React child.}
  and rendered nothing. I learned that
  \texttt{ReactNode} does not include functions.
\item
  \textbf{\texttt{key} is not a prop.} I read
  \texttt{props.key} and got
  \texttt{undefined}, plus a warning saying
  \texttt{key} is not a prop. React removes it before
  the component runs.
\item
  \textbf{Optional rating.} I annotated
  \texttt{\{ rating: number \}} although
  \texttt{Film} has
  \texttt{rating?: number}. The error was
  \texttt{Type 'Film' is not assignable to type '\{ rating: number; \}'}.
  I kept the \texttt{?} or added a default.
\item
  \textbf{Array index.}
  \texttt{const [firstFilm] = films; firstFilm.title}
  gave
  \texttt{error TS18048: 'firstFilm' is possibly 'undefined'.}
  I used \texttt{?.}, an \texttt{if}
  check, or a default.
\end{enumerate}

\end{document}

